\documentclass[11pt]{amsart}

\usepackage{amsmath,amssymb,amsthm,enumitem,hyperref}
\usepackage[margin=1.05in]{geometry}
\usepackage[T1]{fontenc}
\usepackage{lmodern}
\usepackage{microtype}

\hypersetup{colorlinks=true, linkcolor=blue, citecolor=blue, urlcolor=blue}

\newtheorem{theorem}{Theorem}
\newtheorem{lemma}[theorem]{Lemma}
\newtheorem{proposition}[theorem]{Proposition}
\newtheorem{corollary}[theorem]{Corollary}
\theoremstyle{definition}
\newtheorem*{definition}{Definition}
\newtheorem*{question}{Question}
\newtheorem*{example}{Example}
\theoremstyle{remark}
\newtheorem{remark}{Remark}[section]

\newcommand{\C}{\mathcal C}
\newcommand{\D}{\mathcal D}
\newcommand{\B}{\mathcal B}
\newcommand{\R}{\mathcal R}
\newcommand{\F}{\mathcal F}
\newcommand{\res}{\operatorname{res}}
\newcommand{\wt}{\operatorname{wt}}
\newcommand{\mult}{\operatorname{mult}}
\newcommand{\len}{\ell}


\begin{document}


A partition is a weakly decreasing finite list of nonnegative integers.  The \emph{weight} of a partition is the sum of its parts.

\begin{definition}[The Andrews--Dhar partition families]\label{def:families}
Suppose that $n$ is a positive integer.

\smallskip
\noindent
(b) Let $\D_3(n)$ be the set of partitions $\mu$ of $n$ into nonnegative parts such that
\begin{enumerate}[label=(\roman*)]
\item the smallest part of $\mu$ occurs exactly three times;
\item every part strictly larger than the smallest part occurs at most twice.
\end{enumerate}
Let $D_3(n):=|\D_3(n)|$.
For later use in generating functions, we also set
\[
        \D_3(0):=\{(0,0,0)\},\qquad D_3(0):=1.
\]
\end{definition}

For $\mu\in\D_3(n)$, we define
\[
   \tau(\mu):=\#\{\text{parts of }\mu\text{ strictly larger than its smallest part}\}.
\]
For $i\in\{0,1,2\}$, set
\[
   \D_3^{(i)}(n):=\{\mu\in\D_3(n):\tau(\mu)\equiv i\pmod3\},
   \qquad
   D_3^{(i)}(n):=|\D_3^{(i)}(n)|.
\]

The last operation is a smallest-part raising map
\[
        \mathcal T:\R(N)\longrightarrow \D_3^{(0)}(N+1),
\]
where $\R(N)$ is the set of positive partitions $\sigma$ of $N$ satisfying the following two conditions:
\begin{enumerate}[label=(\roman*),leftmargin=2.6em]
\item no part of $\sigma$ occurs three or more times;
\item either $\ell(\sigma)\equiv2\pmod3$, or $\ell(\sigma)\equiv0\pmod3$ and the smallest part of $\sigma$ is unique.
\end{enumerate}
The second alternative presupposes that $\sigma$ is nonempty; in particular, this convention gives $\R(0)=\varnothing$.
For $\sigma\in\R(N)$, define $\mathcal T(\sigma)$ as follows.  If $\ell(\sigma)\equiv2\pmod3$, adjoin a new part $1$.  If $\ell(\sigma)\equiv0\pmod3$, raise the unique smallest part by $1$.  Call the resulting positive partition $\eta$.  If the smallest part of $\eta$ occurs exactly three times, set $\mathcal T(\sigma)=\eta$.  Otherwise, we set
\[
        \mathcal T(\sigma)=\eta\cup(0,0,0).
\]
Therefore, the only zeros ever created by the forward map are exactly three zeros, and they are created only in the final step.

The final step of the bijection changes a length congruence into a smallest-part condition.  The inverse is simple, but the proof below records why zeros occur only as an exact triple.

The map $\mathcal T$ was defined in the introduction.  Its inverse lowers one smallest part.

For $n\ge1$ and $\mu\in\D_3^{(0)}(n)$, define $L(\mu)$ as follows.  If the smallest part of $\mu$ is $0$, delete the three zero parts; if the smallest part is positive, do nothing.  Call the resulting positive partition $\eta$.  Let $s$ be the smallest part of $\eta$.  Lower one copy of $s$ to $s-1$, delete that new part if it is $0$, and sort.  The result is $L(\mu)$.

The following lemma proves both well-definedness and invertibility of the raising step.

\begin{lemma}\label{lem:L}
For $n\ge1$, the maps
\[
        L:\D_3^{(0)}(n)\longrightarrow\R(n-1),
        \qquad
        \mathcal T:\R(n-1)\longrightarrow\D_3^{(0)}(n)
\]
are inverse bijections.
\end{lemma}

\begin{proof}
Let $\mu\in\D_3^{(0)}(n)$.  If the smallest part of $\mu$ is positive, then before lowering, the number of positive parts is $3+\tau(\mu)$.  If the smallest part is $0$, then after deleting the three zeros, the number of positive parts is $\tau(\mu)$.  In both cases this length is divisible by $3$, because $\tau(\mu)\equiv0\pmod3$.

Lowering one smallest positive part lowers the weight by $1$.  If the original smallest part was positive, then three copies of it existed before lowering; after lowering, only two copies remain at that old size, and the newly created smaller part is either deleted because it is zero or occurs once.  All larger parts already occurred at most twice.  If the original smallest part was zero, then, after the three zeros are deleted, all remaining positive parts already occur at most twice.  Hence no part of $L(\mu)$ occurs three or more times.

There are two length cases.  If the lowered part was $1$, then it became $0$ and was deleted; the length changes from $0\pmod3$ to $2\pmod3$.  If the lowered part was greater than $1$, the length remains $0\pmod3$, and the new smallest part is unique.  Therefore $L(\mu)\in\R(n-1)$.

Now start with $\sigma\in\R(n-1)$.  If $\ell(\sigma)\equiv2\pmod3$, the map $\mathcal T$ adjoins a $1$, reversing exactly the case in which a lowered $1$ had been deleted.  If $\ell(\sigma)\equiv0\pmod3$, then the smallest part is unique by definition of $\R$, and $\mathcal T$ raises that part by $1$, reversing exactly the case in which a positive smallest part had been lowered.  After this reversal, either the smallest positive part occurs exactly three times, or it does not.  Apart from the part just adjoined or raised, all positive multiplicities are unchanged from $\sigma$ and hence remain at most two; the only possible creation of a triple positive part is therefore at the new smallest positive part.  If that part occurs exactly three times, the output is already in $\D_3^{(0)}(n)$.  If it does not, the only possible original $D_3$-partition had smallest part $0$, so $\mathcal T$ appends exactly three zeros.  In both cases the number of parts above the smallest part is divisible by $3$, so the output lies in $\D_3^{(0)}(n)$.

The two descriptions are exact reversals, so $L$ and $\mathcal T$ are inverse bijections.
\end{proof}


\end{document}
